% preamble.tex -- packages and macros of the SSCDO#2 porting report.
% Macros are defined here only; check.py parses \misura, \misurau, \studio,
% \scheda, \prova, \questione by name and arity.
\usepackage{fontspec}
\usepackage{unicode-math}
\setmainfont{Libertinus Serif}
\setsansfont{Libertinus Sans}
\setmonofont[Scale=MatchLowercase]{Libertinus Mono}
\setmathfont{Libertinus Math}
\usepackage[italian]{babel}
\usepackage[a4paper,margin=2.3cm,marginparwidth=1.8cm,marginparsep=3mm]{geometry}
\usepackage{amsmath,booktabs,longtable,array,tabularx,enumitem}
\usepackage{siunitx}
\sisetup{output-decimal-marker={,},group-separator={\,},group-minimum-digits=5,
  retain-explicit-plus=true,range-phrase={--},per-mode=symbol}
\usepackage{xcolor}
\definecolor{sdeciso}{HTML}{2E7D32}
\definecolor{sdaprovare}{HTML}{EF6C00}
\definecolor{sdomanda}{HTML}{1565C0}
\usepackage{tikz}
\usetikzlibrary{positioning,arrows.meta}
\usepackage{listings}
\lstdefinestyle{faust}{basicstyle=\ttfamily\small,breaklines=true,
  keepspaces=true,columns=fullflexible,frame=single,framesep=4pt}
\lstset{style=faust}
\usepackage{tcolorbox}
\tcbuselibrary{breakable}
\usepackage[hidelinks]{hyperref}

% ---- state tags
\newcommand{\statotag}[2]{\textcolor{#1}{\textsf{\small\bfseries #2}}}
\newcommand{\deciso}{\statotag{sdeciso}{DECISO}}
\newcommand{\daprovare}{\statotag{sdaprovare}{DA PROVARE}}
\newcommand{\domanda}{\statotag{sdomanda}{DOMANDA}}

% ---- numbers and sources
\newcommand{\misura}[2]{\num{#1}}
\newcommand{\misurau}[3]{\qty{#1}{#2}}
\newcommand{\studio}[1]{\texttt{doc/study/sscdo2/#1/}}
\newcommand{\file}[1]{\texttt{\detokenize{#1}}}
\newcommand{\vuoto}[1]{\ifx\relax#1\relax ---\else #1\fi}

% ---- a card per control
\newcounter{scheda}
\newcommand{\scheda}[8]{%
  \refstepcounter{scheda}\label{scheda:#1}%
  \begin{tcolorbox}[breakable,colback=white,colframe=black!40,boxrule=0.4pt,
    title={\textsf{\bfseries Scheda \thescheda\quad #2}\hfill #7},coltitle=black,colbacktitle=black!6]
  \begin{tabularx}{\linewidth}{@{}>{\bfseries\small\raggedright\arraybackslash}p{3.9cm}X@{}}
  Ambito & \vuoto{#3}\\
  Valore iniziale & \vuoto{#4}\\
  Conseguenze & \vuoto{#5}\\
  Rispetto all'originale & \vuoto{#6}\\
  Studio e audio & \vuoto{#8}\\
  \end{tabularx}
  \end{tcolorbox}}

% ---- a rehearsal card, with blank fields to fill by hand
\newcounter{prova}
\newcommand{\campo}[1]{\par\noindent\textbf{#1:}\ \hrulefill\par\vspace{2mm}\noindent\hrulefill\par}
\newcommand{\prova}[4]{%
  \refstepcounter{prova}\label{prova:#1}%
  \begin{tcolorbox}[breakable,colback=white,colframe=sdaprovare!60,boxrule=0.4pt,
    title={\textsf{\bfseries Prova \theprova}},coltitle=black,colbacktitle=sdaprovare!8]
  \textbf{Ascoltare:} #2\par
  \textbf{Valore:} #3\par
  \textbf{Riferimento:} #4\par\medskip
  \campo{Esito}\campo{Decisione}
  \end{tcolorbox}}

% ---- a numbered question for Davide
\newcounter{questione}
\newcommand{\questione}[2]{%
  \refstepcounter{questione}\label{q:#1}%
  \par\noindent\domanda\ \textbf{Q\thequestione.}\ #2\par\medskip}

% ---- let long monospace names break lines instead of overflowing
\setlength{\emergencystretch}{2em}
